The Requirements Analysis and Specification Document (RASD) is a key reference for planning and estimating the size, cost, and timeline of a project.
It explains what the system is meant to achieve, covering its scope, objectives, and goals.
Along with outlining non-functional requirements, it breaks down functional requirements using use cases, diagrams, and class models.
The RASD also helps with testing, verification, and validation, ensuring the system works as intended.
Moreover, it keeps track of changes and the progress of the system as it evolves over time.

\section{Purpose}\label{sec:purpose}

This document focuses on the Requirements Analysis and Specification Document (RASD) for the Students\&Companies (S\&C) platform.
It contains a detailed description of the platform's primary goals, domain, and relevant models used to represent the domain.
It includes an analysis of fundamental scenarios and their respective use cases, with a list of functional and non-functional requirements that define the development of the software.

Moreover, this document examines interface requirements, performance attributes and constraints that are relevant regarding the platform's user experience.
It also includes a formal Analysis using Alloy which consists in a brief presentation of the main objectives driving the formal modeling activity and the results of the checks performed on meaningful assertions.
For better context, the document includes a history of its development, references, and an explanation of its structure.

The purpose of this document is to guide the creation of the Students\&Companies software, a web application used to match students with companies offering internship opportunities.
\\\\\\\\

%chiedi PER I PACKAGE
\subsection{Goals}\label{subsec:goals}
\renewcommand{\arraystretch}{1.5} % Adjust row height for better readability
% goal counter
\newcounter{g}
\setcounter{g}{1}
\newcommand{\cg}{\theg\stepcounter{g}}
\begin{center}
    \begin{longtable}{|c|p{0.90\textwidth}|}
        \hline
        \rowcolor[HTML]{E7E0FA}
        \textbf{ID} & \textbf{Description} \\ \hline
        \rowcolor[HTML]{FFFFFF}
        \textbf{G\cg} & Allow students and companies to provide information (internships, CVs) \\ \hline
        \rowcolor[HTML]{F0F0F0}
        \textbf{G\cg} & Allow companies to advertise internships \\ \hline
        \rowcolor[HTML]{FFFFFF}
        \textbf{G\cg} & Allow students to search for internships \\ \hline
        \rowcolor[HTML]{F0F0F0}
        \textbf{G\cg} & Allow students and companies to deal with the recommendation process\\ \hline
        \rowcolor[HTML]{FFFFFF}
        \textbf{G\cg} & Allow the management of interview and selections \\ \hline
        \rowcolor[HTML]{F0F0F0}
        \textbf{G\cg} & Allow the Analysis Tool to collect various kinds of information regarding the internships \\ \hline
        \rowcolor[HTML]{FFFFFF}
        \textbf{G\cg} & Allow suggestions to students and companies regarding how to make submissions \\ \hline
        \rowcolor[HTML]{F0F0F0}
        \textbf{G\cg} & Allow students and companies to establish a contact \\ \hline
        \rowcolor[HTML]{FFFFFF}
        \textbf{G\cg} & Allow students and companies to perform the selection process \\ \hline
        \rowcolor[HTML]{F0F0F0}
        \textbf{G\cg} & Allow students and companies to provide feedbacks \\ \hline
        \rowcolor[HTML]{FFFFFF}
        \textbf{G\cg} & Allow universities to monitor internships (complaints) \\ \hline
    \end{longtable}
\end{center}


\section{Scope}\label{sec:scope}
The Students\&Companies platform is designed to ease the matching process between students looking for internships and companies offering them.
S\&C matches students with internships based on skills, experience, and career goals as outlined in students' CVs and company project descriptions.
The platform offers to the students to look proactively for internship but also allows notification mechanisms to alert students when relevant internships become available.
Likewise, companies are notified when a student’s profile aligns with their internship descriptions.
This system employs mechanisms of various level of sophistication like keyword matching and advanced analyses considering students and internship characteristics via external Analysis Tool.

Once an initial contact has been established, the platform supports the selection process, allowing companies to conduct interviews, manage questionnaires, and assess candidates’ compatibility with the internship.

It supports both students and companies by suggesting optimizations for student CVs and company project descriptions to make their submissions more appealing for their counterparts.

S\&C provides interested parties with mechanisms to track internship progress and report any issues.
While universities, after authentication on the platform could monitor these reports to ensure the performance of their students.

The platform ensures secure authentication for companies, students and universities that are affiliated with S\&C.

\subsection{World phenomena}\label{subsec:world-phenomena}
\renewcommand{\arraystretch}{1.5} % Adjust row height for better readability
\begin{center}
    \begin{longtable}{|c|p{0.90\textwidth}|}
        \hline
        \rowcolor[HTML]{D6EAF8} % Light beige
        \textbf{ID} & \textbf{Description} \\ \hline
        \rowcolor[HTML]{FFFFFF} % Light beige
        \textbf{W1} & Students fulfill their CVs \\ \hline
        \rowcolor[HTML]{F0F0F0} % White
        \textbf{W2} & Companies open an internship position \\ \hline
        \rowcolor[HTML]{FFFFFF} % Light beige
        \textbf{W3} & Universities decide to terminate an internship / Universities handle complaints \\ \hline
        \rowcolor[HTML]{F0F0F0} % White
        \textbf{W4} & Students and companies conduct external interviews  \\ \hline
    \end{longtable}
\end{center}

\subsection{Shared phenomena}\label{subsec:shared-phenomena}
\subsubsection{World Controlled}
Phenomena controlled by the world and observed by the machine
\renewcommand{\arraystretch}{1.5} % Adjust row height for better readability
\begin{center}
    \begin{longtable}{|c|p{0.90\textwidth}|}
        \hline
        \rowcolor[HTML]{FFE5B4}
        \textbf{ID} & \textbf{Description} \\ \hline
        \rowcolor[HTML]{FFFFFF} % Light beige
        \textbf{S1} & Users signup/login \\ \hline
        \rowcolor[HTML]{F0F0F0} % White
        \textbf{S2} & Students upload their CVs \\ \hline
        \rowcolor[HTML]{FFFFFF} % Light beige
        \textbf{S3} & Companies upload the internship position \\ \hline
        \rowcolor[HTML]{F0F0F0} % White
        \textbf{S4} & Students proactively look for internships \\ \hline
        \rowcolor[HTML]{FFFFFF} % Light beige
        \textbf{S5} & Matched students and companies accept/reject offers/applications \\ \hline
        \rowcolor[HTML]{F0F0F0} % White
        \textbf{S6} & Companies setup the interview through the platform \\ \hline
        \rowcolor[HTML]{FFFFFF} % Light beige
        \textbf{S7} & Companies collect answers from interviewed students \\ \hline
        \rowcolor[HTML]{F0F0F0} % White
        \textbf{S8} & Companies file a complaint \\ \hline
        \rowcolor[HTML]{FFFFFF} % Light beige
        \textbf{S9} & Students file a complaint \\ \hline
    \end{longtable}
\end{center}



\subsubsection{Machine Controlled}
Phenomena controlled by the machine and observed by the world
\renewcommand{\arraystretch}{1.5} % Adjust row height for better readability
\begin{center}
    \begin{longtable}{|c|p{0.90\textwidth}|}
        \hline
        \rowcolor[HTML]{BDFCC9}
        \textbf{ID} & \textbf{Description} \\ \hline
        \rowcolor[HTML]{FFFFFF} % White
        \textbf{S10} & The system sends complaints to universities \\ \hline
        \rowcolor[HTML]{F0F0F0} % Light beige
        \textbf{S11} & The system asks students to provide feedback and suggestions to feed statistical analysis \\ \hline
        \rowcolor[HTML]{FFFFFF} % White
        \textbf{S12} & The system asks companies to provide feedback and suggestions \\ \hline
        \rowcolor[HTML]{F0F0F0} % Light beige
        \textbf{S13} & The system informs students when a position becomes available \\ \hline
        \rowcolor[HTML]{FFFFFF} % White
        \textbf{S14} & The system informs companies of the availability of student CVs corresponding to their needs \\ \hline
        \rowcolor[HTML]{F0F0F0} % Light beige
        \textbf{S15} & The system helps manage the interviews \\ \hline
        \rowcolor[HTML]{FFFFFF} % White
        \textbf{S16} & The system finalizes the selections \\ \hline
        \rowcolor[HTML]{F0F0F0} % Light beige
        \textbf{S17} & The system provides suggestions regarding the submissions of companies and students \\ \hline
        \rowcolor[HTML]{FFFFFF} % White
        \textbf{S18} & The system provides interested parties with mechanisms to track and monitor the matchmaking process \\ \hline
        \rowcolor[HTML]{F0F0F0} % Light beige
        \textbf{S19} & The system notifies the appropriate students' universities of complaints to handle \\ \hline
    \end{longtable}
\end{center}

\section{Glossary}\label{sec:definitions-acronyms-abbreviations}
\subsection{Definitions}\label{subsec:definition}
\begin{itemize}
    \item \textbf{Student:} A registered user on the S\&C platform who is searching for internship opportunities and uploads their CV.
    \item \textbf{Company:} An entity registered on the platform that publishes internship listings and selects candidates.
    \item \textbf{University:} An organization that supervises internships and handles complaints.
    \item \textbf{CV:} Curriculum Vitae, a document summarizing a student's experiences, skills, and qualifications.
    \item \textbf{Internship:} A time-limited work opportunity offered by a company to train students on specific projects.
    \item \textbf{Recommendation Engine:} A system that analyzes student profiles and internship listings to suggest relevant matches.
    \item \textbf{Complaint:} A report submitted by a student or a company to address an issue related to an internship.
    \item \textbf{Feedback:} Retrospective information provided by students and companies to evaluate the internship experience.
    \item \textbf{User:} Student and companies
    \item \textbf{User*:} Student, companies and universities
    \item \textbf{Guest:} Student, companies and universities not yet registered
\end{itemize}

\subsection{Acronyms}\label{subsec:acronyms}
\begin{itemize}
    \item \textbf{S\&C}: Students\&Companies
    \item \textbf{CV}: Curriculum Vitae
    \item \textbf{API}: Application Programming Interface
    \item \textbf{GDPR}: General Data Protection Regulation
    \item \textbf{UML}: Unified Modeling Language
    \item \textbf{RASD}: Requirements Analysis and Specification Document
    \item \textbf{DD}: Design Document
    \item \textbf{HR}: Human Resource
\end{itemize}

\subsection{Abbreviations}\label{subsec:abbreviations}
\begin{itemize}
    \item \textbf{e.g.}: Example (\textit{exempli gratia})
    \item \textbf{etc.}: And so on (\textit{et cetera})
    \item \textbf{i.e.}: That is (\textit{id est})
    \item \textbf{vs.}: Against (\textit{versus})
    \item \textbf{w.r.t.}: With respect to (\textit{with respect to})
    \item \textbf{ID}: Unique Identifier
    \item \textbf{G}: Goal
    \item \textbf{WP}: World Phenomena
    \item \textbf{SP}: Shared Phenomena
    \item \textbf{R}: Requirement
\end{itemize}

\section{Revision History}\label{sec:revision-history}
\begin{table}[H]
    \centering
    \renewcommand{\arraystretch}{1.5}
    \begin{tabular}{|c|c|p{0.6\textwidth}|}
        \hline
        \rowcolor[HTML]{5C8A99}
        \textbf{Version} & \textbf{Date} & \textbf{Description} \\ \hline
        1.0              & 30/10/2024    & Initial version of the Requirements Analysis and Specification Document. \\ \hline
    \end{tabular}
    \caption{Revision History}
\end{table}

\section{Reference Documents}\label{sec:reference-documents}
\begin{itemize}
    \item Assignment RDD AY 2024-2025: Problem statement and project details.
    \item IEEE Std 830-1998: IEEE Recommended Practice for Software Requirements Specifications.
    \item GDPR (General Data Protection Regulation): European regulation for data protection and privacy.
    \item S\&C Design Document: To be developed in the second phase of the project.
    \item UML 2.0 Standard: Unified Modeling Language for diagrams and models.
    \item Software Engineering 2 Course Materials: Lectures and supplementary notes.
\end{itemize}

This document is organized as follows:
\begin{itemize}
    \item \textbf{Section 1 - Introduction:} Provides an overview of the project, its scope, goals, and relevant background information. It also includes definitions, acronyms, abbreviations, reference documents, and a description of the document structure.
    \item \textbf{Section 2 - Overall Description:} Describes the product's context, key functionalities, user characteristics, and assumptions, dependencies, and constraints. It includes domain models such as class and state diagrams.
    \item \textbf{Section 3 - Specific Requirements:} Details the external interface requirements, functional requirements (use case and sequence diagrams), performance requirements, and design constraints. This section also covers attributes such as reliability, availability, and security.
    \item \textbf{Section 4 - Formal Analysis Using Alloy:} Explains the formal modeling approach, objectives, and results obtained through Alloy validation, including example worlds and assertion checks.
    \item \textbf{Section 5 - Effort Spent:} Summarizes the amount of time each team member has contributed to the creation of this document.
    \item \textbf{Section 6 - References:} Lists all documents, standards, and resources referenced during the development of this document.
\end{itemize}

